\chapter{2022年天津卷}

\section{单选题}

\begin{problem}
设全集 \(U=\{-2,-1,0,1,2\}\)，集合 \(A=\{0,1,2\}\)，\(B=\{-1,2\}\)，则 \(A\cap\left(C_UB\right)=\)（\quad）

\choices
    {\(\{0,1\}\)}
    {\(\{0,1,2\}\)}
    {\(\{-1,1,2\}\)}
    {\(\{-1,0,1,2\}\)}
\end{problem}

\begin{answer}
A
\end{answer}

\begin{solution}
由全集 \(U\) 可得 \(C_UB=\{-2,0,1\}\)，因此
\(A\cap(C_UB)=\{0,1,2\}\cap\{-2,0,1\}=\{0,1\}\)，故选 A.
\end{solution}

\begin{problem}
“\(x\) 为整数”是“\(2x+1\) 为整数”的（\quad）

\choices
    {充分不必要条件}
    {必要不充分条件}
    {充要条件}
    {既不充分也不必要条件}
\end{problem}

\begin{answer}
A
\end{answer}

\begin{solution}
若 \(x\) 为整数，则 \(2x+1\) 是整数，所以“\(x\) 为整数”是“\(2x+1\) 为整数”的充分条件. 反之，取 \(x=\dfrac12\)，则 \(2x+1=2\) 为整数，而 \(x\) 不是整数，所以必要性不成立. 故选 A.
\end{solution}

\begin{problem}
函数 \(y=\dfrac{|x^2-1|}{x}\) 的图象大致为（\quad）

\choices
    {\choicebitmap[width=0.15\paperwidth]{img/2022/tianjin/q3_option_a.png}}
    {\choicebitmap[width=0.15\paperwidth]{img/2022/tianjin/q3_option_b.png}}
    {\choicebitmap[width=0.15\paperwidth]{img/2022/tianjin/q3_option_c.png}}
    {\choicebitmap[width=0.15\paperwidth]{img/2022/tianjin/q3_option_d.png}}
\end{problem}

\begin{answer}
A
\end{answer}

\begin{solution}
函数的定义域为 \(x\ne0\)，且
\[
y=\begin{cases}
x-\dfrac{1}{x}, & x\le -1\text{ 或 }x\ge1,\\
\dfrac{1}{x}-x, & -1<x<0\text{ 或 }0<x<1.
\end{cases}
\]
当 \(x>0\) 时，函数值为正；当 \(x<0\) 时，函数值为负，且 \(f(-x)=-f(x)\). 又因为 \(f(1)=0\)，当 \(0<x<1\) 时函数从正无穷递减到 \(0\)，当 \(x>1\) 时函数从 \(0\) 递增. 因此图象应为选项 A，故选 A.
\end{solution}

\begin{problem}
将 1916 年到 2015 年的全球年平均气温（单位：\({}^{\circ}\mathrm C\)），共 100 个数据，分成 6 组：\([13.55,13.75)\)，\([13.75,13.95)\)，\([13.95,14.15)\)，\([14.15,14.35)\)，\([14.35,14.55)\)，\([14.55,14.75]\)，并整理得到如下的频率分布直方图，则全球年平均气温在区间 \([14.35,14.75]\) 内的有（\quad）

\bitmapfigure[max width=7.2cm]{img/2022/tianjin/q4_fig1.png}

\choices
    {\(22\text{ 年}\)}
    {\(23\text{ 年}\)}
    {\(25\text{ 年}\)}
    {\(35\text{ 年}\)}
\end{problem}

\begin{answer}
B
\end{answer}

\begin{solution}
由直方图可读出各组的组距均为 \(0.20\)，相应的频率分布直方图高度依次为 \(0.40\)，\(1.30\)，\(1.55\)，\(0.60\)，\(0.50\)，\(0.65\).

区间 \([14.35,14.75]\) 对应第五、六组，其频率为
\(0.20(0.50+0.65)=0.23\)，故数据个数为 \(100\times0.23=23\)，选 B.
\end{solution}

\begin{problem}
设 \(a=2^{0.7}\)，\(b=\left(\dfrac{1}{3}\right)^{0.7}\)，\(c=\log_2\dfrac{1}{3}\)，则 \(a\)，\(b\)，\(c\) 的大小关系为（\quad）

\choices
    {\(a<b<c\)}
    {\(c<a<b\)}
    {\(b<c<a\)}
    {\(c<b<a\)}
\end{problem}

\begin{answer}
D
\end{answer}

\begin{solution}
因为 \(a=2^{0.7}>1\)，\(0<b=\left(\dfrac{1}{3}\right)^{0.7}<1\)，而 \(c=\log_2\dfrac{1}{3}<0\)，所以 \(c<b<a\)，故选 D.
\end{solution}

\begin{problem}
化简 \(\left(2\log_4 3+\log_8 3\right)\left(\log_3 2+\log_9 2\right)=\)（\quad）

\choices
    {1}
    {\(\dfrac{5}{4}\)}
    {2}
    {\(\dfrac{5}{2}\)}
\end{problem}

\begin{answer}
C
\end{answer}

\begin{solution}
设 \(t=\log_2 3\)，则 \(\log_4 3=\dfrac{t}{2}\)，\(\log_8 3=\dfrac{t}{3}\)，\(\log_3 2=\dfrac{1}{t}\)，\(\log_9 2=\dfrac{1}{2t}\). 因此原式等于
\[
\left(t+\dfrac{t}{3}\right)\left(\dfrac{1}{t}+\dfrac{1}{2t}\right)=\dfrac{4t}{3}\cdot\dfrac{3}{2t}=2.
\]
故选 C.
\end{solution}

\begin{problem}
已知双曲线 \(\dfrac{x^2}{a^2}-\dfrac{y^2}{b^2}=1\)（\(a>0,b>1\)）的左，右焦点分别为 \(F_1\)，\(F_2\)，抛物线 \(y^2=4\sqrt{5}x\) 的准线经过 \(F_1\)，且与双曲线的一条渐近线交于点 \(A\). 若 \(\angle F_1F_2A=\dfrac{\pi}{4}\)，则双曲线的方程为（\quad）

\choices
    {\(\dfrac{x^2}{16}-\dfrac{y^2}{4}=1\)}
    {\(\dfrac{x^2}{4}-\dfrac{y^2}{16}=1\)}
    {\(\dfrac{x^2}{4}-y^2=1\)}
    {\(x^2-\dfrac{y^2}{4}=1\)}
\end{problem}

\begin{answer}
D
\end{answer}

\begin{solution}
设双曲线的半焦距为 \(c\). 抛物线的准线为 \(x=-\sqrt{5}\)，故 \(F_1=(-\sqrt{5},0)\)，从而 \(c=\sqrt{5}\)，即 \(a^2+b^2=5\). 双曲线的渐近线为 \(y=\pm\dfrac{b}{a}x\)，所以直线 \(x=-\sqrt{5}\) 与渐近线的交点可取 \(A=(-\sqrt{5},\pm\dfrac{b\sqrt{5}}{a})\).

此时 \(F_2=(\sqrt{5},0)\)，故 \(\overrightarrow{F_2F_1}=(-2\sqrt{5},0)\)，且 \(\overrightarrow{F_2A}=(-2\sqrt{5},\pm\dfrac{b\sqrt{5}}{a})\). 由夹角为 \(\dfrac{\pi}{4}\)，得
\[
\tan\angle F_1F_2A=\dfrac{b\sqrt{5}/a}{2\sqrt{5}}=\dfrac{b}{2a}=1,
\]
所以 \(b=2a\). 代入 \(a^2+b^2=5\)，得 \(5a^2=5\)，即 \(a=1\)，\(b=2\). 故双曲线方程为 \(x^2-\dfrac{y^2}{4}=1\)，选 D.
\end{solution}

\begin{problem}
“十字歇山顶”是中国古代建筑屋顶的经典样式之一，图 1 中的故宫角楼的顶部即为十字歇山顶. 其上部可视为由两个相同的直三棱柱交叠而成的几何体（图 2）. 这两个直三棱柱有一个公共侧面 \(ABCD\)，在底面 \(BCE\) 中，若 \(BE=CE=3\)，\(\angle BEC=120^{\circ}\)，则该几何体的体积为（\quad）

\examfiguregroup[float=inline]{img/2022/tianjin/q8_fig1.png}{%
    \begin{tabular}{cc}
    \bitmapinclude[max width=6.0cm]{img/2022/tianjin/q8_fig1.png} &
    \bitmapinclude[max width=6.0cm]{img/2022/tianjin/q8_fig2.png}
    \end{tabular}%
}

\choices
    {\(\dfrac{27}{2}\)}
    {\(\dfrac{27\sqrt{3}}{2}\)}
    {27}
    {\(27\sqrt{3}\)}
\end{problem}

\begin{answer}
C
\end{answer}

\begin{solution}
记底面三角形 \(BCE\) 的面积为 \(S\)，则
\[
S=\frac12\cdot BE\cdot CE\cdot\sin\angle BEC
=\frac12\cdot3\cdot3\cdot\sin120^\circ
=\frac{9\sqrt3}{4}.
\]
由余弦定理，
\[
BC^2=BE^2-2BE\cdot CE\cos120^\circ+CE^2
=9+9+9=27,
\]
所以 \(BC=3\sqrt3\).

由图中两个全等棱柱的对应关系，共同侧面 \(ABCD\) 的边 \(BC\) 在一个棱柱中是底面边，在另一个棱柱中是侧棱；同理，\(CD\) 在前一个棱柱中是侧棱，在后一个棱柱中是底面边. 因此 \(CD=BC\). 令 \(h=BC=CD=3\sqrt3\)，则 \(ABCD\) 为边长 \(h\) 的正方形.
建立空间直角坐标系，使
\(B=(0,0,0)\)，\(C=(h,0,0)\)，\(A=(0,h,0)\)，\(D=(h,h,0)\). 由等腰三角形 \(BCE\) 的对称性，取
\(E=\left(\dfrac h2,0,\dfrac{h}{2\sqrt3}\right)\). 设另一个棱柱底面中与 \(E\) 对应的顶点为 \(G\)，则可取
\(G=\left(0,\dfrac h2,\dfrac{h}{2\sqrt3}\right)\).

在 \(0\leq t\leq h\) 上令
\(\varphi(t)=\dfrac1{\sqrt3}\min\{t,h-t\}\). 第一个棱柱在点 \((x,y,z)\) 上的范围为 \(0\leq z\leq\varphi(x)\)，第二个棱柱的范围为 \(0\leq z\leq\varphi(y)\). 第一个棱柱的体积为 \(Sh\).

考虑第二个棱柱在第一个棱柱外的部分. 当 \(0\leq x\leq y\leq\dfrac h2\) 时，该部分为一个三棱锥，其四个顶点为
\(B\)、\(M=\left(0,\dfrac h2,0\right)\)、\(G\) 以及
\(H=\left(\dfrac h2,\dfrac h2,\dfrac{h}{2\sqrt3}\right)\). 其底面 \(\triangle MGH\) 在平面 \(y=\dfrac h2\) 内，且

\[
S_{\triangle MGH}=\frac12\cdot\frac h2\cdot\frac{h}{2\sqrt3}=\frac{h^2}{8\sqrt3},
\]

顶点 \(B\) 到该底面的距离为 \(\dfrac h2\)，所以该三棱锥的体积为

\[
V_1=\frac13\cdot\frac{h^2}{8\sqrt3}\cdot\frac h2=\frac{h^3}{48\sqrt3}.
\]

由关于 \(x=\dfrac h2\) 和 \(y=\dfrac h2\) 的对称性，第二个棱柱在第一个棱柱外的部分由 4 个全等的上述三棱锥组成. 因此

\[
V=Sh+4V_1
=\frac{9\sqrt3}{4}\cdot3\sqrt3+4\cdot\frac{(3\sqrt3)^3}{48\sqrt3}
=\frac{81}{4}+\frac{27}{4}=27.
\]
故选 C.
\end{solution}

\begin{problem}
关于函数 \(f(x)=\dfrac{1}{2}\sin 2x\)，给出下列结论：

\begin{circlelist}
\item \(f(x)\) 的最小正周期是 \(2\pi\)；
\item \(f(x)\) 在区间 \(\left[-\dfrac{\pi}{4},\dfrac{\pi}{4}\right]\) 上单调递增；
\item 当 \(x\in\left[-\dfrac{\pi}{6},\dfrac{\pi}{3}\right]\) 时，\(f(x)\) 的取值范围为 \(\left[-\dfrac{\sqrt{3}}{4},\dfrac{\sqrt{3}}{4}\right]\)；
\item \(f(x)\) 的图象可以由函数 \(g(x)=\dfrac{1}{2}\sin\left(2x+\dfrac{\pi}{4}\right)\) 的图象向左平移 \(\dfrac{\pi}{8}\) 个单位长度得到.
\end{circlelist}
其中正确结论的个数为（\quad）

\choices
    {1}
    {2}
    {3}
    {4}
\end{problem}

\begin{answer}
A
\end{answer}

\begin{solution}
函数 \(f(x)=\dfrac12\sin2x\) 的最小正周期为 \(\pi\)，所以结论 1 错误.

在区间 \(\left[-\dfrac\pi4,\dfrac\pi4\right]\) 上，\(2x\in\left[-\dfrac\pi2,\dfrac\pi2\right]\)，正弦函数在该区间上单调递增，所以结论 2 正确.

当 \(x\in\left[-\dfrac\pi6,\dfrac\pi3\right]\) 时，\(2x\in\left[-\dfrac\pi3,\dfrac{2\pi}3\right]\). 在此区间上，\(\sin2x\) 的最小值为 \(-\dfrac{\sqrt3}{2}\)，最大值为 \(1\)，故
\(f(x)\in\left[-\dfrac{\sqrt3}{4},\dfrac12\right]\)，不是题给范围，结论 3 错误.

又
\(g\left(x-\dfrac\pi8\right)=\dfrac12\sin\left(2x-\dfrac\pi4+\dfrac\pi4\right)=f(x)\)，所以 \(f(x)\) 的图象由 \(g(x)\) 的图象向右平移 \(\dfrac\pi8\) 个单位长度得到，结论 4 错误. 正确结论只有 1 个，故选 A.
\end{solution}

\section{填空题}

\begin{problem}
已知 \(\i\) 是虚数单位，化简 \(\dfrac{11-3\i}{1+2\i}\) 的结果为\fillinblank{}.
\end{problem}

\begin{answer}
\(1-5\i\)
\end{answer}

\begin{solution}
分子、分母同乘分母的共轭复数 \(1-2\i\)，得
\[
\frac{11-3\i}{1+2\i}
=\frac{(11-3\i)(1-2\i)}{(1+2\i)(1-2\i)}
=\frac{11-22\i-3\i+6\i^2}{5}
=\frac{5-25\i}{5}=1-5\i.
\]
\end{solution}

\begin{problem}
在 \(\left(\sqrt{x}+\dfrac{3}{x^2}\right)^5\) 的展开式中，常数项是\fillinblank{}.
\end{problem}

\begin{answer}
15
\end{answer}

\begin{solution}
展开式的通项为
\[
\mathrm C_5^r(\sqrt{x})^r\left(\frac{3}{x^2}\right)^{5-r}
=\mathrm C_5^r3^{5-r}x^{\frac r2-2(5-r)}
=\mathrm C_5^r3^{5-r}x^{\frac{5r}{2}-10}.
\]
要得到常数项，必须有 \(\dfrac{5r}{2}-10=0\)，即 \(r=4\). 所以常数项为 \(\mathrm C_5^4\cdot3=15\).
\end{solution}

\begin{problem}
若直线 \(x-y+m=0\)（\(m>0\)）被圆 \((x-1)^2+(y-1)^2=3\) 截得的弦长等于 \(m\)，则 \(m\) 的值为\fillinblank{}.
\end{problem}

\begin{answer}
2
\end{answer}

\begin{solution}
圆心为 \(O(1,1)\)，半径为 \(r=\sqrt3\). 圆心到直线的距离为
\[
d=\frac{|1-1+m|}{\sqrt{1^2+(-1)^2}}=\frac{m}{\sqrt2}.
\]
由弦长公式，
\[
m=2\sqrt{r^2-d^2}=2\sqrt{3-\frac{m^2}{2}}.
\]
两边平方得 \(m^2=12-2m^2\)，所以 \(m^2=4\). 又 \(m>0\)，故 \(m=2\).
\end{solution}

\begin{problem}
现有 52 张扑克牌（去掉大小王），每次取一张，取后不放回，则两次都抽到 \(A\) 的概率为\fillinblank{}；在第一次抽到 \(A\) 的条件下，第二次也抽到 \(A\) 的概率为\fillinblank{}.
\end{problem}

\begin{answer}
\(\dfrac{1}{221}\)；\(\dfrac{1}{17}\)
\end{answer}

\begin{solution}
52 张牌中有 4 张 \(A\). 两次都抽到 \(A\) 的概率为
\[
\frac4{52}\cdot\frac3{51}=\frac1{221}.
\]
在第一次抽到 \(A\) 的条件下，剩余 51 张牌中有 3 张 \(A\)，所以第二次也抽到 \(A\) 的条件概率为
\[
\frac3{51}=\frac1{17}.
\]
\end{solution}

\begin{problem}
在 \(\triangle ABC\) 中，点 \(D\) 为 \(AC\) 的中点，点 \(E\) 满足 \(\overrightarrow{CB}=2\overrightarrow{BE}\). 记 \(\overrightarrow{CA}=\bs{a}\)，\(\overrightarrow{CB}=\bs{b}\)，用 \(\bs{a}\)，\(\bs{b}\) 表示 \(\overrightarrow{DE}=\)\fillinblank{}；若 \(\overrightarrow{AB}\perp\overrightarrow{DE}\)，则 \(\angle ACB\) 的最大值为\fillinblank{}.
\end{problem}

\begin{answer}
\(\dfrac{3}{2}\bs{b}-\dfrac{1}{2}\bs{a}\)；\(\dfrac{\pi}{6}\)
\end{answer}

\begin{solution}
以 \(C\) 为起点，用位置向量表示各点，则
\(\overrightarrow{CA}=\bs a\)，\(\overrightarrow{CB}=\bs b\)，且 \(D\) 为 \(AC\) 的中点，所以
\(\overrightarrow{CD}=\dfrac12\bs a\). 由
\(\overrightarrow{CB}=2\overrightarrow{BE}\)，得
\(\overrightarrow{BE}=\dfrac12\bs b\)，从而
\(\overrightarrow{CE}=\overrightarrow{CB}+\overrightarrow{BE}=\dfrac32\bs b\). 因此
\[
\overrightarrow{DE}=\overrightarrow{CE}-\overrightarrow{CD}
=\frac32\bs b-\frac12\bs a.
\]

又 \(\overrightarrow{AB}=\bs b-\bs a\)，所以垂直条件给出
\[
(\bs b-\bs a)\cdot(3\bs b-\bs a)=0,
\]
即
\[
3|\bs b|^2+|\bs a|^2-4\bs a\cdot\bs b=0.
\]
设 \(\theta=\angle ACB\)，并令 \(t=\dfrac{|\bs a|}{|\bs b|}>0\). 由 \(\bs a\cdot\bs b=|\bs a||\bs b|\cos\theta\)，上式化为
\[
t^2-4t\cos\theta+3=0.
\]
因此 \(4t\cos\theta=t^2+3>0\)，从而 \(\cos\theta>0\). 该关于 \(t\) 的方程有正实根，判别式必须满足
\[
16\cos^2\theta-12\geq0.
\]
结合 \(\cos\theta>0\) 和 \(0<\theta<\pi\)，得
\(\cos\theta\geq\dfrac{\sqrt3}{2}\)，所以 \(\theta\leq\dfrac\pi6\). 当 \(\cos\theta=\dfrac{\sqrt3}{2}\)、\(t=\sqrt3\) 时等号可以取得，故最大值为 \(\dfrac\pi6\).
\end{solution}

\begin{problem}
设 \(a\in\R\)，对任意实数 \(x\)，记 \(f(x)\) 表示 \(|x|-2\)，\(x^2-ax+3a-5\) 中的较小者. 若函数 \(f(x)\) 至少有 3 个零点，则 \(a\) 的取值范围为\fillinblank{}.
\end{problem}

\begin{answer}
\([10,+\infty)\)
\end{answer}

\begin{solution}
令 \(p(x)=|x|-2\)，\(q(x)=x^2-ax+3a-5\). 因为 \(f(x)=\min\{p(x),q(x)\}\)，所以 \(f(x)=0\) 当且仅当 \(p(x)\geq0\)，\(q(x)\geq0\)，且 \(p(x)=0\) 或 \(q(x)=0\).

当 \(p(x)=0\) 时，\(x=\pm2\). 分别计算
\(q(2)=a-1\)，\(q(-2)=5a-1\). 因此 \(x=2\) 是零点的条件为 \(a\geq1\)，\(x=-2\) 是零点的条件为 \(a\geq\dfrac15\).

其余零点必须是方程 \(q(x)=0\) 且满足 \(|x|>2\) 的根. 方程 \(q(x)=0\) 的判别式为
\[
\Delta=a^2-4(3a-5)=(a-2)(a-10).
\]
当 \(a<2\) 时，二次函数的顶点横坐标为 \(\dfrac a2<1\)，所以在 \(x\geq2\) 上严格递增，方程 \(q(x)=0\) 在 \(x>2\) 上至多有一个根；若有这样的根，则 \(q(2)<0\)，从而 \(x=2\) 不是 \(f\) 的零点. 在 \(x\leq-2\) 上，若 \(a\geq-4\)，则 \(q'(x)=2x-a\leq-4-a\leq0\)，方程 \(q(x)=0\) 在 \(x<-2\) 上至多有一个根；若 \(a<-4\)，则顶点在 \(-2\) 的左侧，且 \(q(-2)=5a-1<0\)，所以左支至多有一个根落在 \(x<-2\)，另一个根若存在则大于 \(-2\). 因此 \(q(x)=0\) 在 \(|x|>2\) 上至多有两个根，且一旦有 \(x<-2\) 的根，就有 \(q(-2)<0\)，从而 \(x=-2\) 也不是 \(f\) 的零点. 结合 \(x=\pm2\) 的情况可知，此时至多得到两个零点，不能满足至少 3 个零点的要求.

当 \(a=2\) 时，\(q(x)=(x-1)^2\)，其唯一零点 \(x=1\) 不满足 \(p(x)\geq0\)，所以仍只有 \(x=\pm2\) 两个零点. 当 \(2<a<10\) 时，\(\Delta<0\)，故 \(q(x)>0\) 对任意 \(x\) 成立，只有 \(x=\pm2\) 两个零点. \(a=10\) 时，
\(q(x)=x^2-10x+25=(x-5)^2\)，所以 \(x=-2\)，\(x=2\) 和 \(x=5\) 均为零点，恰有 3 个零点.

当 \(a>10\) 时，方程 \(q(x)=0\) 有两个根
\[
x_{1,2}=\frac{a\mp\sqrt{(a-2)(a-10)}}2.
\]
其中较小根满足 \(x_1>2\)：因为 \(x_1>2\) 等价于 \(a-4>\sqrt{(a-2)(a-10)}\)，而 \(a>10\) 时两边均为正，平方后等价于 \(4(a-1)>0\). 故两个根都大于 2，再加上 \(x=\pm2\)，零点数至少为 4.

综上，函数至少有 3 个零点当且仅当 \(a\geq10\)，故填 \([10,+\infty)\).
\end{solution}

\section{解答题}

\begin{problem}
在 \(\triangle ABC\) 中，角 \(A\)，\(B\)，\(C\) 所对的边分别为 \(a\)，\(b\)，\(c\). 已知 \(a=\sqrt{6}\)，\(b=2c\)，\(\cos A=-\dfrac{1}{4}\).

\begin{enumerate}
\item 求 \(c\) 的值；
\item 求 \(\sin B\) 的值；
\item 求 \(\sin(2A-B)\) 的值.
\end{enumerate}
\end{problem}

\begin{solution}
\begin{enumerate}
\item 由余弦定理，\(a^2=b^2+c^2-2bc\cos A\). 代入 \(a=\sqrt6\)，\(b=2c\)，\(\cos A=-\dfrac14\)，得 \(6=4c^2+c^2+c^2=6c^2\). 由于边长 \(c>0\)，所以 \(c=1\).
\item 由 \(a=\sqrt6\) 及 \(\cos A=-\dfrac14\)，得 \(\sin A=\sqrt{1-\cos^2A}=\dfrac{\sqrt{15}}4\). 由正弦定理，\(\dfrac{\sin B}{b}=\dfrac{\sin A}{a}\)，又 \(b=2\)，所以 \(\sin B=\dfrac{2\cdot\sqrt{15}/4}{\sqrt6}=\dfrac{\sqrt{10}}4\).
\item 由余弦定理，\(\cos B=\dfrac{a^2+c^2-b^2}{2ac}=\dfrac{6+1-4}{2\sqrt6}=\dfrac{\sqrt6}{4}\). 又 \(\sin 2A=2\sin A\cos A=-\dfrac{\sqrt{15}}8\)，\(\cos 2A=2\cos^2A-1=-\dfrac78\)，故
\[
\sin(2A-B)=\sin2A\cos B-\cos2A\sin B
=-\dfrac{\sqrt{15}}8\cdot\dfrac{\sqrt6}{4}+\dfrac78\cdot\dfrac{\sqrt{10}}4
=\dfrac{\sqrt{10}}8.
\]
\end{enumerate}
\end{solution}

\begin{problem}
如图，在直三棱柱 \(ABC-A_1B_1C_1\) 中，\(AC\perp AB\)，点 \(D\)，\(E\)，\(F\) 分别为 \(A_1B_1\)，\(AA_1\)，\(CD\) 的中点，\(AB=AC=AA_1=2\).

\begin{enumerate}
\item 求证：\(EF\parallel\) 平面 \(ABC\)；
\item 求直线 \(BE\) 与平面 \(CC_1D\) 所成角的正弦值；
\item 求平面 \(A_1CD\) 与平面 \(CC_1D\) 所成二面角的余弦值.
\end{enumerate}

\bitmapfigure[max width=4.8cm]{img/2022/tianjin/q17_fig1.png}
\end{problem}

\begin{solution}
\begin{enumerate}
\item 建立空间直角坐标系，使 \(A\) 为原点，\(AB\)，\(AC\)，\(AA_1\) 分别为 \(x\) 轴，\(y\) 轴，\(z\) 轴，则 \(A(0,0,0)\)，\(B(2,0,0)\)，\(C(0,2,0)\)，\(A_1(0,0,2)\)，\(B_1(2,0,2)\)，\(C_1(0,2,2)\). 由中点坐标公式得 \(D(1,0,2)\)，\(E(0,0,1)\)，\(F(\dfrac12,1,1)\)，从而 \(\overrightarrow{EF}=(\dfrac12,1,0)\). 该向量的 \(z\) 坐标为 \(0\)，且 \(\overrightarrow{EF}=\dfrac14\overrightarrow{AB}+\dfrac12\overrightarrow{AC}\)，所以 \(EF\parallel\) 平面 \(ABC\).
\item 由 \(\overrightarrow{CC_1}=(0,0,2)\)，\(\overrightarrow{CD}=(1,-2,2)\)，可取平面 \(CC_1D\) 的一个法向量为 \(\bs n=(2,1,0)\). 又 \(\overrightarrow{BE}=(-2,0,1)\)，所以直线 \(BE\) 与平面 \(CC_1D\) 所成角的正弦值为
\[
\dfrac{|\overrightarrow{BE}\cdot\bs n|}{|\overrightarrow{BE}|\cdot|\bs n|}=\dfrac{4}{\sqrt{5}\sqrt{5}}=\dfrac45.
\]
\item 平面 \(A_1CD\) 的一个法向量可取 \(\bs m=(0,1,1)\)，平面 \(CC_1D\) 的一个法向量为 \(\bs n=(2,1,0)\). 因此两平面所成二面角的余弦值为
\[
\dfrac{|\bs m\cdot\bs n|}{|\bs m|\cdot|\bs n|}=\dfrac{1}{\sqrt{2}\sqrt{5}}=\dfrac{\sqrt{10}}{10}.
\]
\end{enumerate}
\end{solution}

\begin{problem}
设 \(\{a_n\}\) 为等差数列，\(\{b_n\}\) 为等比数列，且 \(a_1=b_1=a_2-b_2=a_3-b_3=1\).

\begin{enumerate}
\item 求 \(\{a_n\}\) 和 \(\{b_n\}\) 的通项公式；
\item 记 \(\{a_n\}\) 的前 \(n\) 项和为 \(S_n\)，求证：\(\left(S_{n+1}+a_{n+1}\right)b_n=S_{n+1}b_{n+1}-S_nb_n\)；
\item 求 \(\displaystyle\sum_{k=1}^{2n}\left[a_{k+1}-(-1)^ka_k\right]b_k\).
\end{enumerate}
\end{problem}

\begin{solution}
\begin{enumerate}
\item 设等差数列 \(\{a_n\}\) 的公差为 \(d\)，等比数列 \(\{b_n\}\) 的公比为 \(q\). 由题意，\(a_2=1+d\)，\(a_3=1+2d\)，\(b_2=q\)，\(b_3=q^2\)，所以 \(d=q\)，且 \(2q-q^2=0\). 等比数列的公比不为 \(0\)，故 \(q=2\)，从而 \(d=2\). 因此 \(a_n=1+2(n-1)=2n-1\)，\(b_n=2^{n-1}\).
\item 由 \(a_k=2k-1\) 得 \(S_n=\sum_{k=1}^n(2k-1)=n^2\). 于是
\[
(S_{n+1}+a_{n+1})b_n=((n+1)^2+2n+1)2^{n-1}=(n^2+4n+2)2^{n-1},
\]
而
\[
S_{n+1}b_{n+1}-S_nb_n=((n+1)^2\cdot2^n-n^2\cdot2^{n-1})=(n^2+4n+2)2^{n-1}.
\]
所以 \((S_{n+1}+a_{n+1})b_n=S_{n+1}b_{n+1}-S_nb_n\).
\item 由通项公式，原式为
\[
\sum_{k=1}^{2n}\left[2k+1-(-1)^k(2k-1)\right]2^{k-1}.
\]
当 \(k=2j-1\) 时，对应项为 \((8j-4)4^{j-1}\)；当 \(k=2j\) 时，对应项为 \(4\cdot4^{j-1}\). 因此原式等于
\[
\sum_{j=1}^n8j4^{j-1}=8\sum_{j=1}^n j4^{j-1}.
\]
设 \(T=\sum_{j=1}^n j4^{j-1}\)，则
\[
3T=4T-T=n4^n-\sum_{j=1}^n4^{j-1}=n4^n-\dfrac{4^n-1}{3},
\]
故 \(T=\dfrac{(3n-1)4^n+1}{9}\)，原式为 \(\dfrac{8\left((3n-1)4^n+1\right)}{9}\).
\end{enumerate}
\end{solution}

\begin{problem}
椭圆 \(\dfrac{x^2}{a^2}+\dfrac{y^2}{b^2}=1\)（\(a>b>0\)）的右焦点为 \(F\)，右顶点为 \(A\)，上顶点为 \(B\)，且满足 \(\dfrac{|BF|}{|AB|}=\dfrac{\sqrt{3}}{2}\).

\begin{enumerate}
\item 求椭圆的离心率 \(e\)；
\item 直线 \(l\) 与椭圆有唯一公共点 \(M\)，与 \(y\) 轴相交于 \(N\)（\(N\) 异于 \(M\)）. 记 \(O\) 为坐标原点，若 \(|OM|=|ON|\)，且 \(\triangle OMN\) 的面积为 \(\sqrt{3}\)，求椭圆的标准方程.
\end{enumerate}
\end{problem}

\begin{solution}
\begin{enumerate}
\item 设椭圆的半焦距为 \(c\)，则 \(c^2=a^2-b^2\)，且 \(F=(c,0)\). 因为 \(B=(0,b)\)，所以 \(|BF|=\sqrt{b^2+c^2}=a\)，又 \(|AB|=\sqrt{a^2+b^2}\). 由题设
\[
\dfrac{a}{\sqrt{a^2+b^2}}=\dfrac{\sqrt{3}}{2},
\]
得 \(4a^2=3a^2+3b^2\)，即 \(a^2=3b^2\). 因此 \(e=\dfrac{c}{a}=\sqrt{1-\dfrac{b^2}{a^2}}=\sqrt{\dfrac23}=\dfrac{\sqrt6}{3}\).
\item 设唯一公共点为 \(M=(x_0,y_0)\). 由 \(a^2=3b^2\)，椭圆方程给出 \(x_0^2=3b^2-3y_0^2\). 直线 \(l\) 为椭圆在 \(M\) 处的切线，其方程为 \(\dfrac{x x_0}{a^2}+\dfrac{y y_0}{b^2}=1\)，故它与 \(y\) 轴的交点为 \(N=(0,\dfrac{b^2}{y_0})\). 条件 \(|OM|=|ON|\) 给出
\[
x_0^2+y_0^2=\dfrac{b^4}{y_0^2}.
\]
令 \(t=\dfrac{y_0^2}{b^2}\)，则 \(0<t<1\)，代入 \(x_0^2=3b^2(1-t)\) 得 \(3-2t=\dfrac1t\)，即 \(2t^2-3t+1=0\). 所以 \(t=1\) 或 \(t=\dfrac12\). \(t=1\) 时 \(x_0=0\)，三角形面积为 \(0\)，与题设矛盾，故 \(t=\dfrac12\).

又因为 \(\triangle OMN\) 的面积为 \(\sqrt3\)，所以
\[
\sqrt3=\dfrac12|x_0|\left|\dfrac{b^2}{y_0}\right|=\dfrac{\sqrt3}{2}b^2,
\]
从而 \(b^2=2\)，\(a^2=6\). 故椭圆的标准方程为 \(\dfrac{x^2}{6}+\dfrac{y^2}{2}=1\).
\end{enumerate}
\end{solution}

\begin{problem}
已知 \(a,b\in\R\)，函数 \(f(x)=\e^x-a\sin x\)，\(g(x)=b\sqrt{x}\).

\begin{enumerate}
\item 求函数 \(y=f(x)\) 在点 \((0,f(0))\) 处的切线方程；
\item 若 \(y=f(x)\) 和 \(y=g(x)\) 有公共点，
\begin{enumerate}
\item 当 \(a=0\) 时，求 \(b\) 的取值范围；
\item 求证：\(a^2+b^2>\e\).
\end{enumerate}
\end{enumerate}
\end{problem}

\begin{solution}
\begin{enumerate}
\item \(f(0)=1\)，\(f'(x)=\e^x-a\cos x\)，所以 \(f'(0)=1-a\). 因此曲线在 \((0,f(0))=(0,1)\) 处的切线方程为 \(y-1=(1-a)x\)，即 \(y=(1-a)x+1\).
\item
    \begin{enumerate}
    \item 当 \(a=0\) 时，公共点的横坐标满足 \(x>0\)，并且 \(\e^x=b\sqrt{x}\)，即 \(b=\dfrac{\e^x}{\sqrt{x}}\). 若 \(b\le0\)，等式不可能成立. 令 \(h(x)=\dfrac{\e^x}{\sqrt{x}}\)，则
    \[
    h'(x)=\dfrac{\e^x}{x^{3/2}}\left(x-\dfrac12\right).
    \]
    所以 \(h(x)\) 在 \((0,\dfrac12)\) 上递减，在 \((\dfrac12,+\infty)\) 上递增，最小值为 \(h(\dfrac12)=\sqrt{2\e}\). 故 \(b\in[\sqrt{2\e},+\infty)\).
    \item 公共点的横坐标不能为 \(0\)，设其为 \(x>0\). 由公共点条件，\(\e^x=a\sin x+b\sqrt{x}\). 由柯西不等式，
    \[
    \e^{2x}\le(a^2+b^2)(\sin^2x+x).
    \]
    对任意 \(x>0\)，都有 \(\sin^2x<x\)：当 \(0<x<1\) 时，\(\sin^2x<x^2<x\)；当 \(x\ge1\) 时，\(\sin^2x\le1\le x\)，且 \(x=1\) 时仍为严格不等式. 因此 \(x+\sin^2x<2x\). 又由 \(\e^u\ge1+u\)（取 \(u=2x-1\)），得 \(\e^{2x-1}\ge2x\)，从而 \(\e^{2x-1}>x+\sin^2x\). 于是
    \[
    a^2+b^2\ge\dfrac{\e^{2x}}{x+\sin^2x}>\e,
    \]
    即 \(a^2+b^2>\e\).
    \end{enumerate}
\end{enumerate}
\end{solution}
